\begin{figure}[t]
\centering
\tiny
\begin{tabular}{p{0.47\textwidth} p{0.47\textwidth}}
\multicolumn{1}{c}{\textbf{Raw SFT trace}} & \multicolumn{1}{c}{\textbf{Structured trace (\texttt{<try>/<outcome>})}}\\[3pt]
\begin{tcolorbox}[colback=white, sharp corners, boxrule=0.4pt]
\textbf{Question.} Find all natural numbers that can be written in the form $\frac{ (a+b+c)^{2}}{abc} $, where $ a, b, c $ are natural numbers. \par\medskip
\dots\dots\dots \dots\dots \dots\dots\dots \dots\dots\dots \dots\dots\dots \dots\dots \dots\dots\dots \dots\dots\dots \\
 Wait, how about $a=1,b=1,c=3$? Then $(1+1+3)^2/(1 \cdot 1 \cdot 3) = 5^2/3 = 25/3 \approx 8.333\dots$ which is not a natural number. \dots \dots\dots\dots \dots\dots Wait, so with $a=1, b=4, c=5$, $(1+4+5)^2/(1 \cdot 4 \cdot 5) = 10^2 / 20 = 100/20=5$. So $5$ is possible. Let's check that. $1, 4, 5$: sum is $10$, squared is $100$, divided by $1 \cdot 4 \cdot 5=20$. Yep, $100/20=5$. Perfect. How about $4$? \dots\dots\dots\dots\dots\dots \dots\dots\dots \dots\dots\dots \dots\dots \dots\dots\dots\dots \dots\dots\dots \\

\end{tcolorbox}
&
\begin{tcolorbox}[colback=white, sharp corners, boxrule=0.4pt]
\textbf{Question.} Find all natural numbers that can be written \dots  \par\medskip
% $\frac{ (a+b+c)^{2}}{abc} $, where $ a, b, c $ are natural numbers. \par\medskip
% \dots\dots\dots \dots\dots \dots\dots\dots \dots\dots\dots \dots\dots\dots \dots\dots \dots\dots\dots \dots\dots\dots \\
\trybox{\parbox{1.0\textwidth}{\texttt{<try>} Wait, how about $a=1,b=1,c=3$? Then $(1+1+3)^2/(1 \cdot 1 \cdot 3) = 5^2/3 = 25/3 \approx 8.333\dots$ which is not a natural number. \dots \dots\dots\dots \dots\dots Wait, so with $a=1, b=4, c=5$, $(1+4+5)^2/(1 \cdot 4 \cdot 5) = 10^2 / 20 = 100/20=5$. \dots \dots
% So $5$ is possible. Let's check that. $1, 4, 5$: sum is $10$, squared is $100$, divided by $1 \cdot 4 \cdot 5=20$. 
Yep, $100/20=5$. Perfect. How about $4$? \texttt{</try>}}} \par
\outbox{\parbox{1.0\textwidth}{\texttt{<outcome>} The values $8$ and $6$, are attainable, with $5$ also being attainable through the combination $a=1, b=4, c=5$.
 \texttt{</outcome>}}} \par\medskip
\end{tcolorbox}
\end{tabular}

\caption{\textbf{Raw $\rightarrow$ Structured traces.} We convert free-form solutions into block-structured traces by (i) detecting decision cues (e.g., “Wait” “Hmm”) to \emph{segment} the answer into step-sized spans, (ii) wrapping each span as a \texttt{<try>} (scratch work), and (iii) prompting a larger instruction-tuned model to \emph{summarize} that span into a concise \texttt{<outcome>} (the step’s takeaway). The underlying problem and final answer are unchanged and the trace can be semantically pruned at \texttt{<try>}/\texttt{<outcome>} boundaries.}
\label{fig:raw-to-structured}
\end{figure}


